% Generated by tools/edn_latex.py from reports/edn.md.
% Do not edit: change reports/edn.md and run `make edn`.
\ednentry{
  id = {EDN-71},
  title = {Secrets are scanned by gitleaks before every commit and over the whole history in CI, with a rule of our own for the DagsHub token},
  date = {2026-10-05},
  milestone = {M3: Quality Assurance},
  topic = {CI/CD, Security, Credential handling},
  participants = {@lukas2510, decided by the repository owner},
  decision = {NFR-15 is enforced with gitleaks 8.30.1, run twice with one configuration, \texttt{.\allowbreak{}gitleaks.\allowbreak{}toml}.
The \texttt{gitleaks} pre-commit hook scans the staged change before the commit exists.
The \texttt{Secret scan (gitleaks)} CI job scans every commit GitHub publishes for the repository -- every branch, every tag and the head of every pull request ever opened -- on every pull request and every push to \texttt{main}, with no paths filter, using the release binary pinned by version and SHA-256 rather than \texttt{gitleaks/\allowbreak{}gitleaks-\allowbreak{}action}.
\texttt{.\allowbreak{}gitleaks.\allowbreak{}toml} keeps gitleaks' default rules and adds one, \texttt{dagshub-\allowbreak{}token}, which recognises a DagsHub token by its context because its shape cannot be told from a commit SHA.
A merge commit is scanned for what it changed itself, with \texttt{git log -\allowbreak{}-\allowbreak{}remerge-\allowbreak{}diff}, so a secret written into a conflict resolution is seen as well.
Before the job scans the history it runs \texttt{tools/\allowbreak{}secret\_\allowbreak{}scan.\allowbreak{}py canary}: a throwaway history with a fake token committed in each of thirteen places at its real path, two tokens each, and two merge commits that each bring in one more, scanned with the command and the configuration the real scan uses, which fails unless all 28 are reported for the commit and the file they were planted in.},
  rationale = {\emph{Alternatives considered:}
\begin{itemize}[leftmargin=*, topsep=2pt]
\item \textbf{Option A (chosen): gitleaks, as a pre-commit hook and as a CI job over the whole history.}
\newline Pros: purpose-built, offline and MIT-licensed; the hook stops a secret before it is committed, and the job catches what the hook never saw, a commit made with \texttt{-\allowbreak{}-\allowbreak{}no-\allowbreak{}verify} or from a clone without the hook, including a secret added in one commit and deleted in the next, because it reads the history rather than the tree; the configuration is a file in the repository the course grades, reviewed like code; the default rules stay on, so the credentials the project does not have yet, such as FR-10's API key, a cloud key or a GitHub token, are covered when they arrive.
Everything GitHub publishes for the repository, 228 commits, scans in under two seconds and has no finding, with the default rules or with ours, so no secret was ever committed.
\newline Cons: gitleaks has no rule for a DagsHub token, and its generic rule misses one more often than it looks: in the seven places where the token is assigned with \texttt{=\allowbreak{}} or \texttt{:} it finds 93.0\,\% to 97.5\,\% of them, because 4.9\,\% of such tokens fall below its entropy threshold of 3.5, and in the other six places it finds none, including \texttt{uv run dvc remote modify origin -\allowbreak{}-\allowbreak{}local password <token>}, the exact command the getting-started page documents.
So the project owns a regex, and the canary exists to keep it honest.
The hook is built from source, so its first install is slow, once per machine and gitleaks version: measured on a thermally throttled laptop, 65 s with Go installed and 5 min 35 s without, when pre-commit downloads Go first, at about 145 s of CPU and 300 MB of RAM either way, leaving 47 MB or, with the downloaded Go, 329 MB in pre-commit's cache and about 180 MB in Go's build cache.
Each commit then costs about a quarter of a second of scanning.
\item \textbf{Option B: detect-secrets with a committed baseline.}
\newline Pros: a Python package that installs into our environment, with no Go toolchain; the baseline makes an existing finding auditable rather than a permanent red build.
Its default plugins also detect the DagsHub token better than gitleaks' defaults: in 11 of the 13 places, all 200 tokens in each, missing only the two bare command-line forms, \texttt{dvc remote modify .\allowbreak{}.\allowbreak{}.\allowbreak{} password <token>} and \texttt{dagshub login -\allowbreak{}-\allowbreak{}token <token>}.
\newline Cons: it scans files and, by design, not the git history, so the CI net for a commit that skipped the hook would only see the tip, and a secret added and then deleted would pass both checks; the baseline is a file the team has to keep honest, and an entry audited as ``not a secret'' is an allowlist that nobody re-reviews; and its entropy detectors flag hex and base64 by entropy alone, which on today's tree is three findings to audit -- \texttt{download.\allowbreak{}md5} in \texttt{params.\allowbreak{}yaml}, the gitleaks checksum the CI job pins, and a test string -- and grows with every hash we commit, while the MD5s of \texttt{dvc.\allowbreak{}lock} stay quiet only because it skips lock files.
\item \textbf{Option C: GitHub secret scanning with push protection.}
\newline Pros: nothing to install or maintain, and push protection blocks the push itself rather than reporting it afterwards.
\newline Cons: DagsHub is not among GitHub's supported secret scanning patterns (checked against the documentation and the \texttt{github/\allowbreak{}docs} repository on 2026-10-05), so it would not recognise the one secret the project has; a custom pattern would be another platform setting, as would the feature itself, leaving no trace in the repository the course grades and nothing a reviewer can read in a pull request.
\item \textbf{Option D: a Pytest test that greps the working tree for credential-shaped strings.}
\newline Pros: entirely ours, and the evidence would sit in the test suite, where the matrix could report NFR-15 as verified by a test.
\newline Cons: the weakest detection of the four: it sees only the tree at the commit under test, so a secret that was committed and deleted passes, and it would have to reimplement the context matching every scanner above already has; it would run only in the test job, which has a paths filter and a shallow checkout; and the matrix status it bought would describe a test standing in for the scan rather than the scan.
\end{itemize}
\par
\emph{Why:} What the requirement has to survive is a commit that skipped the hook, and the repository is public, so a secret in any commit of any pushed branch is published whether or not the tip still holds it.
That makes reading the history the property that decides, and it rules out B and D, which read files; C fails earlier, because it cannot see a DagsHub token at all.
Detection is not what decided it, and the measurement says so: detect-secrets' defaults see the token in more places than gitleaks' do.
What decides is that gitleaks' gap can be closed with one rule in a file we own, while B's history gap cannot be closed at all.
Three choices inside A follow from the same evidence.
The rule matches a 40-hex run by its context -- after a word naming a credential or one of the two services, or as the password in a URL -- with word boundaries, so the 32-hex MD5s of \texttt{dvc.\allowbreak{}lock}, \texttt{params.\allowbreak{}yaml} and the \texttt{.\allowbreak{}dvc} files and the 64-hex SHA-256s of \texttt{uv.\allowbreak{}lock} cannot fire it; with it, gitleaks finds all 2,600 planted tokens in all 13 places and still nothing in the tree or the history.
The CI job uses the release binary because \texttt{gitleaks/\allowbreak{}gitleaks-\allowbreak{}action} requires a licence key for a repository owned by an organisation, as this one is, obtained through a registration form and stored as a repository secret, which a pull request from a fork would not receive; the binary is checked against a SHA-256 pinned in the workflow, not one downloaded beside it.
And the job scans every ref, the heads of merged pull requests included, rather than only the pull request's own commits, because the cost is seconds and a leak on someone else's branch, or in a pull request whose branch a squash merge deleted, is published too: it turns every pull request red until it is revoked and acknowledged, which is the point, and a rule added later reaches all of it.
The canary is there because a scan that reports nothing is only evidence if it could have reported something: a regex edited into matching nothing or an allowlist widened over \texttt{.\allowbreak{}env} would otherwise turn the job silently green.
Its first version planted the tokens as files and scanned the directory, and an independent review showed what that cannot see: the real scan reads paths relative to the repository, so an allowlist anchored at the root, \texttt{\textasciicircum{}\textbackslash{}.\allowbreak{}env\$}, hid a committed \texttt{.\allowbreak{}env} token from it while the canary still passed.
The same review found that \texttt{git log -\allowbreak{}p}, which gitleaks runs, shows a merge commit no diff, so a token written into a conflict resolution and committed with \texttt{-\allowbreak{}-\allowbreak{}no-\allowbreak{}verify} was in the published history and in no scanned diff.
\texttt{-\allowbreak{}-\allowbreak{}remerge-\allowbreak{}diff} closes that by showing what a merge changed against an automatic remerge of its parents; \texttt{-\allowbreak{}-\allowbreak{}diff-\allowbreak{}merges=\allowbreak{}first-\allowbreak{}parent} caught it as well, but it reports again, for every merge of \texttt{main} into a branch, everything \texttt{main} ever carried, and each of those would need its own \texttt{.\allowbreak{}gitleaksignore} entry.
The canary now plants a history and scans it with the same command as the real scan, and it was checked against each way of blinding it: with \texttt{.\allowbreak{}gitleaks.\allowbreak{}toml} it finds all 28 tokens; with the default rules alone it misses 20; with the \texttt{\textasciicircum{}\textbackslash{}.\allowbreak{}env\$} allowlist, or a \texttt{.\allowbreak{}gitleaksignore} entry without its commit, it misses 7; and without \texttt{-\allowbreak{}-\allowbreak{}remerge-\allowbreak{}diff} it misses the two merge commits.
The hook is the \texttt{golang} variant because it needs nothing installed by hand, where \texttt{gitleaks-\allowbreak{}system} needs gitleaks on the path and \texttt{gitleaks-\allowbreak{}docker} needs Docker and an unpinned image.
CI skips it in the lint job, because it scans only staged changes, of which CI has none, and would report a scan of nothing as passed.
The job becomes one of the checks branch protection requires once this pull request is merged, which the owner approved; it is a repository setting, so it is not in the diff.
NFR-15 is \textbf{[manual]} in the specification because its evidence is this CI job and not a Pytest test.},
  ai = {Information seeking, Alternative generation, Alternative assessment, Recommendation, Solution generation},
  response = {Accepted with modifications},
  assessment = {AI laid out options A to D with their trade-offs and recommended A, and the owner chose it.
AI then tested the assumptions the choice rested on before implementing it, and one of them did not hold as stated: detect-secrets had been described as weaker on detection and noisy on the MD5s of \texttt{dvc.\allowbreak{}lock}, and measured, its defaults detect the DagsHub token in more places than gitleaks' and it skips \texttt{dvc.\allowbreak{}lock} entirely.
The decision survives on the argument that did hold, history scanning, and this entry records the measured reasons rather than the assumed ones.
The measurement also produced the parts of the design that were not in the recommendation: the \texttt{dagshub-\allowbreak{}token} rule, after the default rules turned out to miss the token in six places outright and in one case in twenty elsewhere, and the canary that keeps that rule from failing silently.
The modification came from an independent AI review of the pull request, which reproduced two holes in that first implementation, the directory-mode canary and the merge commits nothing scanned, and both were closed before merging.},
  aievidence = {Claude Code session on 2026-10-05 implementing issue \#60: the token's shape was read from the gitignored \texttt{.\allowbreak{}env} as length and character classes only, without printing the value; fake tokens of that shape were planted and scanned with both tools (\texttt{reports/\allowbreak{}analysis/\allowbreak{}secret\_\allowbreak{}scan\_\allowbreak{}dagshub\_\allowbreak{}token.\allowbreak{}py} and its results file); gitleaks-action's licence terms were read from its README and the repository's owner type from the GitHub API; the hook's install cost was measured with \texttt{/\allowbreak{}usr/\allowbreak{}bin/\allowbreak{}time} on a clean pre-commit cache, once with Go removed from the path and once with Go 1.27 on it; and a token committed with \texttt{-\allowbreak{}-\allowbreak{}no-\allowbreak{}verify} and deleted in the next commit was confirmed caught by the history scan.
The canary and each blinding variant were run in a container with git 2.45, because the laptop's git 2.34 predates \texttt{-\allowbreak{}-\allowbreak{}remerge-\allowbreak{}diff}; what an older Go does with gitleaks' \texttt{go.\allowbreak{}mod} was checked the same way with Go 1.22.},
  evidence = {\href{https://github.com/mlops-2627q1-mds-upc/MLOPS_Recommenditos/issues/60}{issue \#60}; \href{https://github.com/mlops-2627q1-mds-upc/MLOPS_Recommenditos/pull/70}{PR \#70}; \texttt{.\allowbreak{}gitleaks.\allowbreak{}toml}; \texttt{tools/\allowbreak{}secret\_\allowbreak{}scan.\allowbreak{}py}; the \texttt{secret-\allowbreak{}scan} job in \href{https://github.com/mlops-2627q1-mds-upc/MLOPS_Recommenditos/blob/main/.github/workflows/ci.yml}{\texttt{.\allowbreak{}github/\allowbreak{}workflows/\allowbreak{}ci.\allowbreak{}yml}} and the gitleaks hook in \href{https://github.com/mlops-2627q1-mds-upc/MLOPS_Recommenditos/blob/main/.pre-commit-config.yaml}{\texttt{.\allowbreak{}pre-\allowbreak{}commit-\allowbreak{}config.\allowbreak{}yaml}}; \href{https://github.com/mlops-2627q1-mds-upc/MLOPS_Recommenditos/blob/main/reports/analysis/secret_scan_dagshub_token.py}{\texttt{reports/\allowbreak{}analysis/\allowbreak{}secret\_\allowbreak{}scan\_\allowbreak{}dagshub\_\allowbreak{}token.\allowbreak{}py}} and \href{https://github.com/mlops-2627q1-mds-upc/MLOPS_Recommenditos/blob/main/reports/analysis/secret_scan_dagshub_token_results.txt}{its results}; \href{https://github.com/mlops-2627q1-mds-upc/MLOPS_Recommenditos/blob/main/docs/docs/specification.md}{specification} NFR-15; the Secrets section of \href{https://github.com/mlops-2627q1-mds-upc/MLOPS_Recommenditos/blob/main/CONTRIBUTING.md\#secrets}{CONTRIBUTING.md}; \hyperref[EDN-46]{EDN-46}, \hyperref[EDN-66]{EDN-66}.},
}
